% LSPV-3 REPORT::PHABSORB1::2026-10-03T21:31:00-07:00
\documentclass[11pt]{article}
\usepackage[margin=.85in]{geometry}
\usepackage{amsmath,amssymb}
\begin{document}
\section*{LSPV-3 PHABSORB1: fixed-level PH provenance predicates are absorbed under the P=NP counterfactual}
\textbf{STATUS: OPEN.} P versus NP remains open. The theorem below is \textbf{CONDITIONAL} on \(P=NP\) and awaits independent red-team review before promotion in the living synthesis. It proves no SAT lower bound and no nonrelativizing or nonalgebrizing hinge.

\section*{CORPUS REVIEWED}
This run surveyed the repository tree and read root main.tex, AGENTS.md, CODEX.md, the pnp-research skill, current tasks, the newest root synthesis, materially relevant LSPV-1--5 reports, FULLSKOLEM1, STATICNP1, WRAPINV1, the HEC/CRS-PPC/CCSPC causal branch, governance notes, and the active reconstruction/exponent ledger. The predecessor LSPV-2 directory still has no code-search hit attributable to 2026-10-03; the two-cycle gap is already documented by research/governance/2026-10-03_0621_LSPV2-TWO-CYCLE-WRITE-GAP.tex, so no missing mathematics was reconstructed.

Primary literature checked this run includes Akshay et al., SAT 2026, on polynomial Skolem/Herbrand extraction from a specified QBF proof system; Chew--Peitl, SAT 2026, on dependency schemes and proof checking; and Boehm--Beyersdorff, SAT 2026, on CAQE proof-theoretic strategy extraction. These are model-specific proof/QBF results and do not supply the arbitrary-static-SAT bridge.

\section*{Exact selector theorem}
Fix a SAME\_REALIZATION input \(X\) and polynomially bounded challenge \(u\). Let \(w\in\{0,1\}^{p(|X|)}\) encode every demanded response object, including any proposed polynomial-size provenance/history string, local algebraic opening, source tag, reconstruction witness, or tuple thereof. Let \(V(X,u,w)\) be a fixed-level polynomial-hierarchy predicate. Assume totality:
\[
\forall u\ \exists w\ V(X,u,w).
\]

\paragraph{Theorem PHABSORB1.}
\textbf{CONDITIONAL on \(P=NP\).} Under \(P=NP\), there is a deterministic polynomial-time uniform selector \(F\) such that
\[
\forall u\quad V(X,u,F(X,u)).
\]

\paragraph{Proof.}
Since \(V\in PH\) and \(P=NP\) implies \(PH=P\), the language
\[
L=\{(X,u,p):\exists s\ V(X,u,ps)\}
\]
also lies in \(PH=P\), with \(ps\) ranging over strings of the prescribed polynomial response length. Starting with the empty prefix, query \(L\) for the prefix extended by 0; choose 0 if an accepting completion exists and otherwise choose 1. Totality guarantees an accepting extension at every stage. There are polynomially many bits and each extension query is polynomial time under the collapse. The resulting \(w\) satisfies \(V(X,u,w)\). \(\square\)

\section*{Quantifier consequence}
This is not a quantifier swap. It preserves challenge dependence:
\[
\forall u\ \exists w\ V(X,u,w)
\]
becomes a uniform polynomial-time function \(w=F(X,u)\), not one \(w\) valid for all \(u\).

FULLSKOLEM1 is recovered by taking \(w\) to be the response \(z\) after the outer \(\sigma\) is selected. PHABSORB1 is stronger in the direction needed by SAME\_REALIZATION auditing: adjoining any polynomial-size history \(h\) to the response does not evade the collapse if validity of \((z,h)\) remains a fixed-level PH predicate.

\section*{No-go corollary for provenance-only repairs}
Suppose a repair strengthens semantic reconstruction to \(B(X,u,z,h)\), where \(|z|+|h|\le |X|^{O(1)}\), \(B\in PH\), and
\[
\forall u\exists z,h\ B(X,u,z,h).
\]
Then under \(P=NP\) a valid pair \((z,h)\) is polynomial-time constructible. Therefore adding polynomially bounded provenance, canonical-opening data, a local algebraic certificate, or a source label whose acceptance remains in fixed PH cannot by itself force superpolynomial acquisition.

This does not show every source-matched condition lies in PH. It gives a falsification test: the next hinge must identify exactly which hypothesis fails, or PHABSORB1 absorbs it.

\section*{HEC, CRS-PPC, CCSPC, and TC4 audit}
HEC and CRS-PPC restricted-model inequalities remain valid in their declared models. PHABSORB1 blocks a common staticization attempt: if delayed-audit or causal-response obligations compile into a polynomially balanced fixed-level PH response relation, then under \(P=NP\) the response can be recomputed by a polynomial selector. A successful static-to-online compiler must preserve a resource not reducible to such response search.

For CCSPC, canonicalization removes witness-choice freedom but does not escape the theorem when the canonical response can be characterized and prefix-searched inside fixed PH. For TC4, the missing statement remains universal: every arbitrary static SAT decider must be forced to realize the extra resource, rather than merely being able to produce an accepted response.

\section*{PQ technique ledger}
Construction-specific collision extraction or special soundness remains \textbf{CANDIDATE}. It survives PHABSORB1 only if the load-bearing conclusion is not merely existence of a polynomial-size accepted opening in fixed PH, but an unconditional source-matched invariant forced on every relevant computation. Generic cryptographic binding and random-oracle binding remain \textbf{FALSIFIED/RETIRED} as unconditional transfers. Akshay et al.'s QBF extraction is a \textbf{KNOWN THEOREM} in its proof-system model and a structural analogy here; acquisition of the proof representation is not free.

\section*{Relativization and algebrization audit}
The argument relativizes. For any oracle \(O\),
\[
P^O=NP^O\Longrightarrow PH^O=P^O,
\]
and the same prefix selector works relative to \(O\). Hence PHABSORB1 cannot be the missing nonrelativizing hinge. It also supplies no nonalgebrizing mechanism.

\section*{Hostile controls}
DEAD, COPY, DELAY, INTERP, and RECOMP remain active. RECOMP is decisive: when the acceptance relation is fixed-level PH and total, the selector reconstructs a fresh valid response without preserving preferred history. Horn, 2-SAT, affine/Tseitin, bounded-treewidth, compact decision diagrams, structured circuits, arithmetic representations, and memoized implementations remain mandatory controls.

\section*{DUPLICATION / CONTRADICTION AUDIT}
PHABSORB1 generalizes FULLSKOLEM1 rather than contradicting it. STATICNP1 is the NP-staticization special case. WRAPINV1 is orthogonal and strengthens the universal-algorithm guard. ENVELOPE1 and descendants remain quantitative accounting results for declared routes; PHABSORB1 gives no certified numerical exponent. Generic quantifier commutation remains \textbf{FALSIFIED/RETIRED}. TransportCert remains accounting, not a lower-bound theorem.

\section*{WHAT CHANGED RELATIVE TO ROOT MAIN}
Root main already states that FULLSKOLEM1 defeats superpolynomial semantic Skolem acquisition under \(P=NP\). The new narrowing is that adding any polynomially bounded response-side provenance or certificate fields does not help when the strengthened validity relation remains fixed-level PH and total. The next candidate must state explicitly which PHABSORB1 hypothesis fails.

\section*{ADOPTED / REJECTED / RETIRED PRIOR COMPONENTS}
\textbf{Adopted:} SAME\_REALIZATION discipline; FULLSKOLEM1; STATICNP1; WRAPINV1; quantitative reconstruction ledgers at their stated scopes; HEC/CRS-PPC/CCSPC as restricted-model or candidate bridge objects.

\textbf{Rejected:} polynomial provenance length, canonical response uniqueness, or polynomial certificate verification as sufficient acquisition hardness.

\textbf{FALSIFIED/RETIRED as generic hinges:} fixed-level-PH response acquisition under \(P=NP\); generic canonical binding; generic cryptographic binding; semantic multiplicity arguments.

\section*{STALE OR OVERSTATED CLAIMS FOUND}
Any claim that adding a polynomial-size history/provenance witness to the displayed reconstruction relation can by itself restore superpolynomial acquisition is overstated unless it identifies why the resulting total search relation is not fixed-level PH or why the selector fails the exact SAME\_REALIZATION obligation.

\section*{Next target}
Source-transcribe the strongest proposed provenance/local-algebraic predicate \(B\) for one exact vMS/SGL realization and classify:
\[
\boxed{\text{polynomial balance? fixed-level PH validity? totality? universal SAT-decider necessity?}}
\]
If the first three answers are yes, semantic acquisition is absorbed under \(P=NP\), and only the fourth, a genuinely computation-forcing invariant outside ordinary response search, can remain load-bearing.

\section*{References}
S. Akshay, O. Beyersdorff, S. Chakraborty, L. Kasche, M. Mahajan, and L. N. Spachmann, Proof Systems for QBF Synthesis: Extracting Skolem and Herbrand Functions, SAT 2026, LIPIcs 377, Article 3.

L. Chew and T. Peitl, Strong (D)QBF Dependency Schemes via Pure Paths with Applications to Proof Checking, SAT 2026, LIPIcs 377, Article 11.

B. Boehm and O. Beyersdorff, Towards Understanding the Complexity of CAQE: A Proof-Theoretic Analysis of Its Core Procedure, SAT 2026, LIPIcs 377, Article 7.

\section*{Final status}
PHABSORB1 is \textbf{CONDITIONAL} on \(P=NP\) pending independent red-team review. SAME\_REALIZATION necessity, TC4, source-matched no-sharing, route completeness, and nonrelativizing/nonalgebrizing hinges remain \textbf{OPEN}. P versus NP remains \textbf{OPEN}.

% END_OF_VERIFIED_REPORT::LSPV-3::PHABSORB1
\end{document}
